\section*{Capítulo \ref{ch:b1:lipids} --- Lípidos}

\begin{solution}{exo:b1:lipids:1}
Un \omterm{def:b1:lipids:fattyacid}{lípido} es una molécula biológica insoluble en agua y soluble en
disolventes apolares. Las \omterm{def:b1:lipids:triglyceride}{grasas}, los \omterm{def:b1:lipids:amphiphilic}{fosfolípidos}, los \omterm{def:b1:lipids:amphiphilic}{esteroles} y los
carotenoides no comparten ningún esqueleto común, solo ese comportamiento
frente al agua, que es lo que les da sus papeles comunes (reservas,
películas, membranas).
\end{solution}

\begin{solution}{exo:b1:lipids:2}
C18:2 $\Delta^{9,12}$. El último doble enlace empieza en el carbono 12
contando desde el carboxilo, es decir, en el carbono 6 desde el extremo
metilo: omega-6.
\end{solution}

\begin{solution}{exo:b1:lipids:3}
Un \omterm{def:b1:lipids:triglyceride}{triglicérido} no tiene cabeza \omterm{def:b1:water-small-molecules:water}{polar}: enteramente hidrófobo, el agua lo
excluye como fase separada, una gota. Un \omterm{def:b1:lipids:amphiphilic}{fosfolípido} es anfífilo: su cabeza
\omterm{def:b1:water-small-molecules:water}{polar} debe quedarse en el agua mientras que sus colas deben abandonarla, y
solo una lámina de dos moléculas de espesor satisface ambas cosas.
\end{solution}

\begin{solution}{exo:b1:lipids:4}
Unos \qty{41}{\celsius}, \qty{-5}{\celsius} y \qty{20}{\celsius} (la curva
del \omterm{def:b1:lipids:amphiphilic}{colesterol} no tiene transición brusca; la mitad de su ascenso está cerca
de \qty{20}{\celsius}). Una bacteria a \qty{10}{\celsius} necesita la
composición insaturada, fluida muy por debajo de su temperatura de
crecimiento.
\end{solution}

\begin{solution}{exo:b1:lipids:5}
$15\,000\times 38 = \qty{570}{MJ}$; $570/8 = 71$ días. Glucógeno: seco,
$570\,000/17 = \qty{33.5}{kg}$; hidratado, \qty{134}{kg}.
\end{solution}

\begin{solution}{exo:b1:lipids:6}
C22:0 $>$ C18:0 $>$ C14:0 $>$ C18:1 $>$ C18:3. Entre los ácidos saturados,
las cadenas más largas establecen más contactos de \omterm{def:b1:water-small-molecules:bonds}{van der Waals}; un doble
enlace \emph{cis} quita más empaquetamiento del que aportan cuatro carbonos
de más (el C18:1 funde por debajo del C14:0); cada enlace adicional lo baja
otra vez.
\end{solution}

\begin{solution}{exo:b1:lipids:7}
$2\times\qty{140}{\micro m^2}/\qty{0.6}{nm^2} = 2\times 1.4\times
10^{-10}/6\times 10^{-19} = \num{4.7e8}$ moléculas; masa
$\num{4.7e8}\times 750/\num{6e23} = \qty{5.8e-13}{g} = \qty{0.58}{pg}$.
\end{solution}

\begin{solution}{exo:b1:lipids:8}
Volumen \qty{1.11}{mL} $= \qty{1.11e-6}{m^3}$; superficie de las esferas
$= 6V/d = 6\times 1.11\times 10^{-6}/10^{-6} = \qty{6.7}{m^2}$. Una sola
gota de ese volumen: $d = \qty{1.28}{cm}$, superficie \qty{5.2}{cm^2}: la
emulsión tiene trece mil veces más interfase. La lipasa es hidrosoluble y
solo actúa en la interfase; la velocidad es proporcional a ella.
\end{solution}

\begin{solution}{exo:b1:lipids:9}
Por encima de $T_m$ sus anillos rígidos estorban el movimiento de las
cadenas vecinas y reducen la \omterm{def:b1:lipids:fluidity}{fluidez}; por debajo de $T_m$ su volumen impide
que las cadenas se empaqueten en el gel ordenado e impide la congelación.
Como no es posible ningún empaquetamiento cooperativo, no hay ninguna
temperatura a la que toda la bicapa cambie de estado a la vez: la transición
queda repartida.
\end{solution}

\begin{solution}{exo:b1:lipids:10}
Con todas las cadenas saturadas: a \qty{25}{\celsius} las membranas están ya
cerca de su transición y rígidas; a \qty{5}{\celsius} gelifican. Las
membranas gelificadas hacen que las proteínas dejen de funcionar y, al
recalentarse o al sufrir un esfuerzo mecánico, se agrietan y se derraman:
las \omterm{def:b1:cell-unit-of-life:cell}{células} pierden su contenido y el \omterm{def:b1:body-plans-tissues:tissue}{tejido} muere (daño por frío). La
planta silvestre desatura sus \omterm{def:b1:lipids:fattyacid}{lípidos} en otoño y sigue fluida a
\qty{5}{\celsius}.
\end{solution}

\begin{solution}{exo:b1:lipids:11}
\qty{1.07}{kg} de agua por kilogramo de \omterm{def:b1:lipids:triglyceride}{grasa}. Oxígeno: \qty{2000}{L}; con
una extracción del \qty{5}{\%} de un \qty{21}{\%} de oxígeno, el aire
necesario es $2000/(0.21\times 0.05) = \qty{190}{m^3}$; agua perdida a
\qty{30}{g/m^3}: \qty{5.7}{kg}. El camello pierde respirando cinco veces más
agua de la que rinde la \omterm{def:b1:lipids:triglyceride}{grasa}: la joroba es una reserva de energía, no de
agua.
\end{solution}

\begin{solution}{exo:b1:lipids:12}
Dada una bicapa, el \omterm{prop:b1:water-small-molecules:properties}{efecto hidrófobo} hace que se vuelva a sellar, que se
cierre en vesículas y que acepte \omterm{def:b1:lipids:fattyacid}{lípidos} nuevos sin aporte de energía: el
ensamblaje es espontáneo. Pero la \omterm{def:b1:cell-unit-of-life:cell}{célula} fabrica sus \omterm{def:b1:lipids:fattyacid}{lípidos} con enzimas
alojadas en una membrana ya existente (el RE), que los insertan allí mismo;
una membrana nueva crece por expansión de una vieja y se reparte en la
división, de modo que toda membrana de toda \omterm{def:b1:cell-unit-of-life:cell}{célula} desciende de las
membranas de la \omterm{def:b1:cell-unit-of-life:cell}{célula} anterior. El autoensamblaje explica la física de la
lámina, no su origen en la \omterm{def:b1:cell-unit-of-life:cell}{célula}.
\end{solution}

\begin{solution}{pb:b1:lipids:1}
\textbf{1.} C16:0 (0), C18:1 $\Delta^9$ (1), C22:6 $\Delta^{4,7,10,13,16,19}$
(6).
\textbf{2.} Aguas cálidas: $0.40\times 0 + 0.45\times 1 + 0.15\times 6 = 1.35$;
aguas frías: $0 + 0.40 + 2.10 = 2.50$ dobles enlaces por cadena.
\textbf{3.} Cada enlace \emph{cis} acoda la cadena e impide el
empaquetamiento estrecho con las vecinas; menos contactos de \omterm{def:b1:water-small-molecules:bonds}{van der Waals},
menos calor necesario para desordenar las cadenas, menor $T_m$.
\textbf{4.} $1.15$ enlaces más por cadena: unos \qty{45}{\celsius} menos.
\textbf{5.} En un gel los \omterm{def:b1:lipids:fattyacid}{lípidos} no pueden moverse: los \omterm{def:b1:membranes-transport:transporters}{transportadores} no
pueden cambiar de conformación, las \omterm{def:b1:membranes-transport:transporters}{bombas} se paran, los \omterm{def:b1:membranes-transport:transporters}{canales} se
atascan; los gradientes iónicos se agotan, la conducción nerviosa falla y el
pez queda paralizado y muere del frío que un pez aclimatado tolera.
\textbf{6.} Las desaturasas, que introducen dobles enlaces en las cadenas de
acilo graso; son \omterm{def:b1:membranes-transport:proteins}{proteínas integrales} del \omterm{def:b1:eukaryotic-cell:endomembrane}{retículo endoplasmático}, donde se
fabrican los \omterm{def:b1:lipids:fattyacid}{lípidos}.
\textbf{7.} En ambos casos tampona la \omterm{def:b1:lipids:fluidity}{fluidez}: endurece la membrana fría,
muy insaturada, por encima de su baja $T_m$, y evita que la membrana cálida
gelifique en un día fresco.
\textbf{8.} $30\,000\times 38 = \qty{1140}{MJ}$.
\textbf{9.} $1140/60 = 19$ días.
\textbf{10.} Seco: $\num{1.14e6}/17 = \qty{67}{kg}$; hidratado,
\qty{268}{kg}.
\textbf{11.} $268 - 30 = \qty{238}{kg}$, casi la mitad de la masa del
camello.
\textbf{12.} Una capa de \omterm{def:b1:lipids:triglyceride}{grasa} bajo toda la piel aislaría el cuerpo y le
impediría evacuar calor en el desierto; concentrar la \omterm{def:b1:lipids:triglyceride}{grasa} en una sola
joroba deja libre el resto de la piel para perder calor.
\textbf{13.} $1.07\times 30 = \qty{32}{kg}$ de agua.
\textbf{14.} $2.0\times 30\,000 = \qty{60000}{L} = \qty{60}{m^3}$ de
oxígeno.
\textbf{15.} $60/(0.21\times 0.05) = \qty{5700}{m^3}$ de aire.
\textbf{16.} $5700\times 30 = \qty{171}{kg}$ de agua.
\textbf{17.} Pierde cinco veces lo que gana: oxidar la joroba cuesta agua.
La joroba es una reserva de energía que permite al camello pasar sin
alimento; su economía de agua procede de sus riñones, de su tolerancia a la
deshidratación y de su capacidad de dejar subir su temperatura.
\textbf{18.} $500\times 3.5\times 6 = \qty{10.5}{MJ}$ almacenados como calor y liberados de noche; evaporarlos habría costado $10\,500/2.4 = \qty{4.4}{kg}$ de agua al día.
\textbf{19.} Una membrana solo puede estirarse un pequeño porcentaje, así
que duplicar el volumen de un glóbulo rojo exige un gran exceso de membrana
plegada en la forma bicóncava de reposo, y un \omterm{def:b1:eukaryotic-cell:cytoskeleton}{citoesqueleto} que le permita
desplegarse sin rasgarse: los glóbulos rojos del camello son ovalados, con
más membrana por unidad de volumen que los de un ser humano.
\textbf{20.} $2\times 80\times 10^{-12}/0.65\times 10^{-18} = \num{2.5e8}$.
\textbf{21.} $\num{2.5e8}\times\num{8e12}\times 40 = \num{7.9e22}$.
\textbf{22.} $\num{7.9e22}\times 750/\num{6e23} = \qty{99}{g}$.
\textbf{23.} La difusión lateral mantiene la cabeza en el agua y las colas
en el \omterm{def:b1:lipids:triglyceride}{aceite} a cada paso, y no cuesta nada; el volteo exige que la cabeza
\omterm{def:b1:water-small-molecules:water}{polar} atraviese el núcleo hidrófobo, una barrera de energía grande que la
agitación térmica supera una vez por semana.
\textbf{24.} Una diferencia entre monocapas decae solo tan deprisa como se
voltean los \omterm{def:b1:lipids:fattyacid}{lípidos}; a una vez por semana y por molécula, una asimetría
establecida por enzimas que trasladan \omterm{def:b1:lipids:fattyacid}{lípidos} concretos (flipasas) persiste
indefinidamente frente al volteo espontáneo.
\textbf{25.} Unos \qty{240}{kg} ahorrados --- casi la mitad de su masa
corporal --- por llevar \qty{30}{kg} de \omterm{def:b1:lipids:triglyceride}{grasa} en vez de \qty{270}{kg} de
glucógeno hidratado; y la joroba no es una reserva de agua: oxidarla pierde
en la respiración cinco veces más agua de la que produce.
\end{solution}
